% =====================================================================
%  CONCLUSIONES  —  un párrafo por objetivo específico, que es una de las
%  dos formas que admite la plantilla y la que menos margen deja a la
%  vaguedad. 2026-10-15: sin códigos del repositorio ni fechas de bitácora.
% =====================================================================
\chapter*{Conclusiones}
\addcontentsline{toc}{chapter}{CONCLUSIONES}

\begin{notaavance}[conclusiones parciales: uno de los cuatro objetivos sigue
abierto]
Estas conclusiones se escriben con el proyecto todavía en curso, y solo para lo
que ya está cerrado. Los objetivos primero, segundo y tercero están verificados,
así que sus conclusiones son definitivas. El cuarto depende de la campaña con los
dos vehículos, que se corre el 16 de octubre de 2026; después de esa fecha no se
toman más datos. La conclusión sobre el objetivo general se enuncia porque
descansa en la campaña de simulación, ya cerrada, y en la navegación de cada
vehículo, ya verificada; falta ejecutar el relevo con los dos vehículos reales.
\end{notaavance}

El primer objetivo específico, modelar la arquitectura funcional del sistema, se
cumplió. El sistema tiene un contrato de interfaces escrito, treinta y seis
requisitos con su prueba, y una capa de coordinación que ejecuta el relevo entre
niveles. Escribir el planificador como función pura, sin depender de la capa de
comunicaciones de ROS~2, resultó más valioso de lo previsto. Permite verificar las
930 combinaciones de origen y destino en milisegundos, y encontró defectos que una
corrida en el simulador no habría delatado, como planes que mandaban al mismo
robot dos veces a la misma pose. La asignación de tareas es óptima por la
estructura del problema, con un solo agente admisible por nivel, y el aporte del
trabajo está en coordinar dos agendas que dependen una de otra durante el relevo.

El segundo objetivo, la plataforma robótica de dos vehículos, se cumplió. Los dos
vehículos están integrados, y Nav2 quedó configurado para cinemática tipo
automóvil, con el planificador, el controlador y el árbol de comportamiento
sustituidos por los que esa cinemática admite (sección~\ref{sec:nav2-ackermann}).
La odometría cumple su criterio, con un error entre $-4{,}6$~\% y $+3{,}3$~\% sobre
recorridos de 6,68 y 14,57~m, y los directores dieron por superada la prueba
correspondiente. La navegación de un vehículo con la llegada verificada también
se superó. En una cadena de tres tramos sin intervención, las llegadas medidas con
flexómetro quedaron entre 0,06 y 0,45~m de la meta, dentro de la tolerancia de
0,5~m (sección~\ref{sec:res-g3}). Los dos vehículos operan juntos en la misma red,
con un tiempo de ida y vuelta de los mensajes de 5,99~ms de mediana.

El límite que queda en la plataforma es la banda muerta de tracción
(sección~\ref{sec:banda-muerta}). El vehículo no avanza por debajo de 0,40~m/s y no
tiene régimen de aproximación fina, y por eso la tolerancia de llegada con los
vehículos es de 0,5~m y no de 0,25~m, como en simulación.

El tercer objetivo, la interfaz de interacción con el usuario, se cumplió. La
interfaz funciona desde el navegador de un teléfono, sin instalación y sin salida
a internet, y se operó con ella una misión completa con relevo desde un
dispositivo real. Dos funciones se añadieron después de ver el sistema en uso.
Una es la confirmación de cambio de piso, que impide dar por completada una
misión con el usuario en el nivel equivocado. La otra es la cancelación, que no
funcionaba porque el servidor de acción rechazaba las solicitudes sin informar.

El cuarto objetivo, la evaluación experimental, está cumplido en simulación y
parcialmente con los vehículos. La campaña de treinta misiones sorteadas fue
válida, sin descartes, y permite afirmar con 95~\% de confianza que la tasa de
éxito supera el 70~\%. Esa parte del objetivo está cerrada y sus cifras no van a
cambiar. Con los vehículos ya se superaron las pruebas de odometría y de
navegación de un vehículo. Falta la campaña con los dos vehículos, que se
concluye al terminar la toma de datos (sección~\ref{sec:res-completo}).

Sobre el objetivo general, el esquema de relevo es viable. Dos robots dedicados
cada uno a un nivel pueden repartirse el guiado de una persona que cambia de piso,
sin que ninguno atraviese la zona de transición vertical y sin modificar el
edificio. Eso quedó demostrado en simulación, con una campaña medida y válida.
Sobre los vehículos reales están verificadas la odometría y la navegación de cada
uno, que durante meses hicieron dudar de que el esquema fuera ejecutable fuera de
la simulación. Falta ejecutar el relevo con los dos vehículos, que es el objeto
de la campaña con los vehículos.

Queda una conclusión de método. En este proyecto costó más averiguar qué medía
cada prueba que construir el sistema, como muestran tres casos. El tiempo de
asignación valía cero en todas las grabaciones, porque el reloj de la simulación
no tenía resolución para medirlo, hasta que se caracterizó en un banco aparte. El
criterio de continuidad entre niveles no podía dar un resultado negativo, así que
sus treinta verificaciones no eran evidencia. Y la correspondencia entre el modelo
del edificio y el edificio real se daba por supuesta, y resultó falsa en un 20~\%
la primera vez que se midió. Los tres se encontraron preguntando si el criterio
podía fallar, y ninguno habría aparecido corriendo más experimentos.
